\section{The idea}

Your idea is that a language model is a store of what words mean. A store of a given size can
hold only so much. So a language with few words should need only a small model, and a model
that is kept the same size should do worse when it is given a language with more words, once
the words no longer fit.

To test this we need a language in which we control the number of words exactly and know every
right answer, so that a model can be graded without doubt. No real language allows that, so we
made one up. In it, ``what a word means'' is something we can write down: a short list of facts
about the thing the word names.

\section{The made-up language}

The language is about animals. Its smallest version has 43 words:

\begin{quote}\small
16 animals: \s{foxes bears wolves owls bats frogs snakes crows}\\
\phantom{16 animals: }\s{hens ducks deer mice toads robins otters lizards}\\
4 colors: \s{red brown green black}\qquad 4 foods: \s{meat plants insects seeds}\\
4 homes: \s{woods fields rivers caves}\qquad 4 waking times: \s{dawn noon dusk night}\\
11 sentence words: \s{are do eat live in wake at yes no . ?}\quad (the full stop and the question mark count as words)
\end{quote}

Each animal has four facts: its color, its food, its home and the time it wakes. Each fact has
one of four values. For example, foxes are red, eat meat, live in woods and wake at dusk. Those
four facts are what the word \s{foxes} means in this language.

There are two kinds of sentence. A \emph{statement} gives one fact. A \emph{question} asks
whether a fact has a certain value and is followed by yes or no:

\begin{quote}\small
\s{foxes are red .}\qquad \s{foxes eat meat .}\qquad \s{foxes live in woods .}\qquad
\s{foxes wake at dusk .}\\
\s{do foxes eat meat ? yes}\qquad \s{do foxes eat seeds ? no}
\end{quote}

Every statement is true and every yes or no is right. The text a model learns from is simply
these sentences one after another, in random order: six in ten are statements and four in ten
are questions. The questions are there so that the language is more than one sentence shape
repeated. The score in this report uses statements only.

\textbf{A larger language} has more animals: 64, 256 or 1,024 in place of 16. The added animals
have invented names (\s{bakos}, \s{gigas}, \s{vofos}, \dots) and invented facts. The facts are
drawn at random, but evenly: a quarter of the animals are red, a quarter brown and so on, and
the same for food, home and waking time. The first 16 animals are the same in every language
and are spread just as evenly. Nothing else differs between the languages: the sentences have
the same shapes and the 27 words that are not animals are the same.

\textbf{Counting how much there is to know.} To say how full a model is, we need a measure of
how much a language asks it to remember. We count in \emph{bits}: one bit is what it takes to
settle a choice between two things, and a choice among four takes two bits. So one fact counts
as 2 bits and an animal, with its four facts, as 8 bits. This is the same for every animal in
every language, as you insisted: a large vocabulary must differ from a small one only in the
number of words, not in how much each word means. A whole language holds 8 bits times its
number of animals: 128 bits for 16 animals, 8,192 for 1,024.

\section{The model and how it learns}\label{sec:model}

A model is a long list of numbers together with a fixed recipe for using them. Given the words
read so far, the recipe produces, for every word of the language, how likely that word is to
come next. The numbers are called \emph{parameters}. At the start they are random, so the
model's predictions are worthless.

\textbf{Training.} The model is given a stretch of text, about 1,160 sentences. At each word it
says how likely every word is to come next, and this is compared with the word that really
comes next. Then all its numbers are changed a little, in the direction that would have made
the real next words more likely. That is one \emph{training step}.

\textbf{The two parts of a model.} The \emph{word table} holds one row of numbers for each
word. The \emph{body} takes the rows of the words read so far and works out the prediction.
The count of numbers in one row is called the \emph{width}: a model of width 8 keeps 8 numbers
for each word. The body is built from pieces of the same width, so a wider model has a larger
body too. Width is the only thing we changed to make a model smaller or larger. We used widths
4, 8, 16 and 32 (table \ref{tab:sizes}). For scale, the models behind chat assistants have
billions of parameters.

\textbf{The same count of numbers for every language, but not the same amount in use.}
Normally a model built for more words would be bigger, because it would need more rows. To
prevent that, every model has a row for each of the 1,024 animals, whichever language it is
trained on, so a model of width 8 always has 10,936 parameters. This only half works. A model
trained on 16 animals never reads the other 1,008 animal words, so their rows lie idle. In
effect it is a smaller model, with its body and 43 rows: about 2,900 numbers at width 8. A model
trained on 1,024 animals uses all 10,936. So every added animal arrives with a row of its own,
8 numbers at width 8, in which its facts could be kept. The only part that is truly the same
size for every language, and that all the animals share, is the body. This matters for reading
the results: if a model learns a larger language without trouble, that may be because each
new animal brought its own row, and not because the model had room to spare. Section
\ref{sec:meaning} comes back to this.

\begin{table}[H]
\centering\small
\begin{tabular}{@{}rrrrcc@{}}
\toprule
Width & \multicolumn{3}{c}{Number of parameters} & Bits of facts per parameter & How full, against \\
 & in the body & in the word table & in all & with all 1,024 animals & a limit of 2 bits \\
\midrule
4 & 880 & 4,204 & 5,084 & 1.61 & about 80\% \\
8 & 2,528 & 8,408 & 10,936 & 0.75 & about 37\% \\
16 & 8,128 & 16,816 & 24,944 & 0.33 & about 16\% \\
32 & 28,544 & 33,632 & 62,176 & 0.13 & about 7\% \\
\bottomrule
\end{tabular}
\caption{The four model sizes. The word table has 1,051 rows in every model: 27 for the words
that are not animals and 1,024 for the animals. The fifth column is the 8,192 bits of the
largest language divided by the model's number of parameters. The last column compares that
with the rough limit of 2 bits per parameter explained below.}
\label{tab:sizes}
\end{table}

\textbf{How full can a model get?} A study of much larger models of this kind (Allen-Zhu and
Li, 2024; a million to several hundred million parameters) found that a model holds at most
about 2 bits of facts per parameter. Nobody has measured this for models as small as ours, so
2 is only a rough guide. By that guide \emph{none} of our models should be full with 8 bits
per animal: width 4 would be about four-fifths full if it learned all 1,024 animals, width 8 a
little over one-third full, and the two widest far emptier. We could see this before the runs.
I ran the plan all the same, because it is the experiment we agreed and because the guide
might not hold for models this small.

\textbf{Three rounds.} We had agreed on three rounds of runs. The \emph{first round} is the
plan as described so far (section \ref{sec:round1}). The \emph{unequal-use round} repeats it
with some animals mentioned far more often than others (section \ref{sec:zipf}). The
\emph{32-bit round} gives each animal four times as much to be known (section
\ref{sec:round32}); it was to be run only if the first round supported your idea.

\section{Keeping the comparison fair}

When the vocabulary grows, nothing else should change, or a worse score could be blamed on
whatever else changed. Two safeguards are already described: every animal carries the same
number of bits, and every model has the same count of numbers whatever the language (though,
as just explained, a larger language puts more of them to use). There are two more.

\textbf{The text is dealt out in equal portions, which we called sets.} A set is always 61,440
sentences, whatever the language. That is what it takes to bring up each of the 4,096 facts of
the largest language 15 times: nine times as a statement, three times in a question answered
yes, and three times in a question answered no. In the first round a set brings up each fact
exactly as often as every other fact, in every language, and each new set is the same sentences
in a new random order. One set is 53 training steps.

\textbf{One more thing changes with the vocabulary, and it cannot be avoided.} A set is the
same length for every language. In the largest language it brings up each fact 15 times. The
smallest language has 64 times fewer facts, so there each fact comes round 960 times in a
set. You pointed out that this is what a larger vocabulary is: each word gets a smaller share
of everything that is said. So that this does not decide the result, the planned comparison is
not made after equal numbers of sets. Each run goes on receiving sets until more training no
longer raises its score, or until its time runs out (next section). Section \ref{sec:equal}
makes the other comparison as well.

\section{The score, and how a run ends}\label{sec:score}

\textbf{The score} is called \emph{facts known}. The model is given the start of a statement,
such as \s{foxes eat}, and nothing else: no earlier sentences it could copy from. If the word
it rates most likely to come next is the right one (\s{meat}), the fact counts as known. This
is done for every fact of every animal in the language. It is a test of memory, on purpose:
the facts were all in the training text, and the question is whether the model has kept them.

A model that knows nothing can still be right one time in four, by always giving the same
value. The score is set up to remove that luck: \textbf{0 stands for one fact in four right,
which is no knowledge at all, and 100 for every fact right.} In between, a score of 90 means
92.5 facts in 100 are right.

A \emph{run} is one model trained on one language. It is graded often at first and then less
and less often: after 1, 2, 3, 4, 6, 8, 11, 16 sets and so on, so that every second grading
comes after about double the training. A run ends in one of three ways, and which way it ends
decides what it can tell us.
\begin{itemize}
\item \textbf{Solved.} The score was 97 or more at two gradings in a row. The model is big
      enough for this language.
\item \textbf{Flat.} The score has stopped rising: doubling the training added less than 2
      points, and doubling it again also added less than 2 points. (An invented example: 70
      after 256 sets, 71 after 512, 71.5 after 1,024.) Before a run is called flat, three more
      things must hold. A second, finer measure, how well the model predicts the next word of
      a text, must have stopped improving too. The score must be at least 5. And the run must
      be past its start: at least 2,000 training steps (about 38 sets in the first round),
      with each fact brought up at least 400 times. The score of a flat run was to be taken as
      the model's limit on that language. Section \ref{sec:round32} shows that this trusted
      the test too much.
\item \textbf{Not finished.} The time allowed ran out first. The time limit is a safety net so
      that one run cannot take the whole night: 30 minutes in the first round, which is a
      little over 1,600 sets. When its time is up a run is graded one last time. If that makes two gradings
      of 97 or more in a row, it counts as solved. Otherwise it is not finished: it tells us
      nothing about the model's limit, and I draw no conclusion about that limit from it.
\end{itemize}
The marks 97 and 2 points here, and the 95 and 10 points of the rule below, are judgment calls
that were fixed before any run. The mark of 5 and the 2,000 steps were added later, but
before the real runs (section \ref{sec:changes}).

\section{The rule we fixed in advance}\label{sec:rule}

Each pairing of a width and a language was to be run three times. I call these the three
\emph{repeats}. They differ in the random starting numbers of the model, in the order of the
sentences, and in the invented facts of the added animals. Within one repeat, models of all
widths read exactly the same text.

Before any run we agreed what would count as clear support for your idea. The rule looks at
one width at a time. In one repeat, that width counts in favor of the idea if all three of
these hold:
\begin{enumerate}[label=(\alph*),leftmargin=2.2em]
\item On the 16-animal language the model scores 95 or more. So it can learn the small language.
\item On one of the larger languages its run ends \emph{flat}, with a score at least 10 points
      below its own score on 16 animals. So it does worse with more words, and not because
      it was stopped early.
\item On that same language, in the same repeat, the model twice as wide scores at least 10
      points higher than this model. So a bigger model wins back what was lost, which points
      to size as the cause.
\end{enumerate}
The idea is clearly supported for that width if this happens in at least two of its three
repeats.

We also agreed how to read the other outcomes. If a model does worse on a larger language but
the model twice as wide does no better, something other than size is holding it back. If no
model does worse on any larger language, the round has not tested the idea: the idea says that
more words hurt once the model is full, and nothing shows that any model was full.

Notice what this means: \textbf{the rule can come out in favor of your idea, or undecided,
but never against it.} We agreed on nothing that would show the idea wrong. The nearest thing
would be a model that ought to be nearly full by the guide of 2 bits per parameter and still
learns every fact, and even that would be weak evidence, because the guide is rough.
